\documentclass[a4paper]{article}

\usepackage{graphicx}
\usepackage{amsmath,amssymb}
\usepackage{minted}
\usepackage{multirow}
\usepackage{float}
\usepackage{todonotes}
\usepackage{forest}
\usepackage{xstring}
\usepackage{xcolor}
\usepackage[most]{tcolorbox}
\usepackage{xparse}
\usepackage{natbib}

\usetikzlibrary{graphs,graphdrawing}
\usegdlibrary{layered}

% for external graphviz
\input{/home/donkarlo/Dropbox/repo/nd_language_ai_project/src/nd_language_ai/formal/computer/programming/tex/latex/code_clips/external_graphviz.tex}

% for tags
\input{/home/donkarlo/Dropbox/repo/nd_language_ai_project/src/nd_language_ai/formal/computer/programming/tex/latex/part/control_sequence/kind/semtag/semtag.tex}

% for links
\input{/home/donkarlo/Dropbox/repo/nd_language_ai_project/src/nd_language_ai/formal/computer/programming/tex/latex/code_clips/seehere.tex}

\title{}
\date{}

\begin{document}
    \maketitle
    
    \section{Definition}
    
    
    The mind--body problem concerns how mental life, including thought and conscious experience, is related to physical states, events, and processes in the body and brain \cite{chambliss-2018-the-mind-body-problem}. It asks whether mental states can be fully explained in physical terms or whether they possess properties that cannot be reduced to physical processes. Different solutions characterize the relationship between mental and physical phenomena in different ways. Physicalist approaches regard mental phenomena as ultimately dependent on or realized by physical states. The problem remains a central issue connecting philosophy of mind, consciousness research, and the cognitive and neural sciences.
    
    The mind–body problem is a philosophical problem concerning the relationship between thought and consciousness in the human mind and body. It addresses the nature of consciousness, mental states, and their relation to the physical brain and nervous system.
    \url{https://en.wikipedia.org/wiki/Mind%E2%80%93body_problem}
    
    \cite{georgiev-2020-quantum-information-theoretic-approach-to-the-mind-brain-problem}  argues that consciousness can be understood as unobservable quantum information integrated in quantum brain states. It distinguishes the private, unobservable mind from the observable brain by using the quantum distinction between states and observables. It also uses Holevo’s theorem to explain why conscious experience is only partially communicable through classical information.
    
    \bibliographystyle{plainnat}
    \bibliography{/home/donkarlo/Dropbox/repo/nd_shared_research/bibliography/bibliography.bib}
\end{document}